<\documentclass[makeidx]{статья}
<<<P000017>{xspace}
\usepackage{epsfig}
\usepackage[utf8]{inputenc}
\usepackage[russian]{babel}
<<<P000019>{xcolor}
\usepackage{syntax}
<<<P000021>[fleqn]{amsmath}
<<<P000022>{amssymb}
<<<P000023>{semantic}
<<<P000024>{dart}}
<<<P000025> [скрыто] [гиперреф]
<<<P000026><lmodern>
<<<P000027>[T1]{fontenc}
\usepackage{makeidx}
<<P000029>
<<<p000030>> Спецификация языка программирования Dart <<<P000036>
[6-е издание]\\
<<<P000031> Версия 2.13-dev
\author{}

% для получения информации о маркерах местоположения (и, в частности,
Команды % <<<P000033> и <<P000034>, см. длинный комментарий на сайте WEB
% конец этого файла.

% изменения
% =========
%
% Существенные изменения в спецификации. Обратите внимание, что версии указаны
% ниже указывают текущую версию цепочки инструментов, когда эти изменения были сделаны.
% На практике новые функции всегда интегрировались в язык.
% спецификация (данный документ) через некоторое время после принятия изменения в
% от используемого языка. По состоянию на сентябрь 2022 года предстоящая версия
% язык, который указывается, указывается номером версии
% скобки после версии цепочки инструментов.
%
% Обратите внимание, что приведенные ниже номера версий (до 2,15) были связаны с
% язык и инструменты, выпущенные в настоящее время на момент изменения спецификации;
% они не были согласованы с фактической версией языка
%, указанных в данной версии документа. Это очень обманчиво, поэтому мы
% изменили субтитры на месяц, а не на номер версии. Точно так же
% «Версия», указанная на первой странице, была изменена для указания
% версия языка, которая будет определена следующей стабильной
% к выпуску данного документа.
%
% к июлю 2025 года
% - Уточнить, что операторы "[]" и "[]=" включены в список "пустых"
% правило о формальных параметрах типа «пустота» и т.д.
%
% май 2025
% - Добавить "s?.length" в качестве постоянного выражения в случае, когда "s" удовлетворяет
% некоторые требования.
%
% апреля 2025 года
% - изменить правила относительно констант, чтобы они были более последовательными: Каждая константа
Выражение % также является потенциально постоянным выражением. Кроме того, устранить
% увольнения в правилах о постоянном сборе букв и постоянных
% выражений объектов (они теперь определены только один раз, вне «Констант»).
%
% Сеп 2024
% - Уточнить правило применимости расширения, чтобы прямо заявить, что
% относится к экземпляру с тем же именем, а не к
% статический член.
%
% Июнь 2024
% - Добавить недостающие ссылки на раздел «Динамика типа» в точках, где
Указывается % статический анализ вызовов членов Объекта.
%
% Февраль 2024
% - Исправить правило о выводе будущего типа из типа (который используется)
% на «плоский».
%
% декабря 2023
% - Разрешить "~/" на операнды типа "двойной" в постоянных выражениях, выравнивая
% спецификации с уже реализованным поведением.
% - Расширить правило грамматики об «инициализаторе» для соответствия
% реализованного поведения. Укажите, что выражение инициализатора не может быть
% - функция буквальная.
% Укажите, в каких ситуациях является ошибкой объявление инициализации
% формальный параметр.
%
% ноябрь 2023
% - Укажите, что динамическая ошибка вызова функции в отложенном и
% еще не загруженная библиотека появится до фактической оценки аргументов;
% после.
%
% Октябрь 2023
% Ввести правило, согласно которому заявление о продлении не может иметь названия
% «тип». Это необходимо для того, чтобы разъяснить "тип расширения".
% называется «на».
%
% Август 2023
% - Исправить текст об ошибке встроенного идентификатора и превратить его в комментарий
% (нормативный текст в правилах грамматики с использованием «typeIdentifier»)
% исправить грамматику для импортных префиксов, чтобы использовать «идентификатор типа».
%
% на июль 2023 года
% - Обновление нулевой безопасности для разделов о переменных.
% - изменить терминологию: «Статическая переменная» — это переменная, объявление которой% включает «статическую» (которую раньше называли «классовой переменной»)
* * * * * * * * * * * * * * * * * * * * * * * * * * * * * * * * * * * * * * * * * * * * * * * * * * * * * * * * * * * * * * * * * * * * * * * * * * * * * * * * * * * * * * * * * * * * * * * * * * * * * * * * * * * * * * * * * * * * * * * * * * * * * * * * * * * * * * * * * * * * * * * * * * * * * * * * * * * * * * * * * * * * * * * * * * * * * * * * * * * * * * * * * * * * * * * * * * * * * * * * * * * * * * * * * * * * * * * * * * * * * * * * * * * * * * * * * * * * * * * * * * * * * * * * * * * * * * * *
%, что это «библиотечная переменная или статическая переменная», где последняя
% новое значение этой фразы. Это позволяет избежать путаницы, которая возникла
% почти каждый раз, когда фраза «статическая переменная» появлялась в
% обсуждение.
% - изменить терминологию: Понятие «изменяемых» и «неизменных» переменных имеет
% исключено (всегда уточняется, какие модификаторы должны быть изменены);
% присутствующих или не присутствующих. Понятие запутывается, когда переменная может
% являются поздними и окончательными, а назначения к ним допускаются статичным анализом.
% - изменить определение «типа элемента функции генератора»
Процентная проблема с текущим определением.
%
% март 2023
% - Уточните, как обрабатываются разрывы линий в многострочной строке в буквальном смысле. имя
% от лексического токена NEWLINE до LINE<<<P000038> BREAK (подтверждает, что это не «P000035»).
% - Очистить правила грамматики (избегать 'T?' как суперкласса/и т.д.
% ошибка; изменить имена расширений на «typeIdentifier»
% встроенных идентификаторов.
%
% Февраль 2023
% - Изменение спецификации постоянных выражений формы "e1 == e2"
%, чтобы использовать примитивное равенство.
%
% декабря 2022
% - Изменить определение функции типа "flatten" для разрешения звука
% выпуск, ср. выпуск SDK #49396.
% - Внесение изменений, связанных с нулевой безопасностью, в разделы о переменных
% и локальных переменных.
% - изменить "примитивный оператор ==" на "примитивное равенство" и включить
% ограничения на «hashCode», соответствующие тем, которые мы имеем на «==».
%
% 2,15 %
% - Позволяет генерическую инстантацию выражений с общим типом функции
% (до сих пор это была ошибка, если выражение не обозначало декларацию).
% - Допустим обобщенную инстанциацию метода «вызова» функционального объекта.
% - Уточнить, что "TypeLiteral.extensionMethod()" является ошибкой, в соответствии с
Погрешность % для «int.toString()».
% - Добавить поддержку клозулизации функций для вызывающих вызов объектов, в том смысле, что
%, что «о» Он называется o.call. Когда тип контекста является типом функции.
% - Уточнить обработку "ковариантных" параметров в интерфейсе класса
%, которые наследуют реализацию, где эти параметры не являются ковариантными.
% - Корректировка и уточнение простой интерполяции строк (чтобы позволить "$это");
Процент уже реализован и полезен.
% - добавьте несколько лексических правил об идентификаторах, уточняющих различные виды.
% - Уточнить конфликты между членами-расширителями и «объектом»
% членов.
% - Правильная <partDeclaration> для включения метаданных.
% - Уточнить раздел о присваиваемых выражениях.
%
% 2,14 %
% - Добавить ограничение по типу параметра, который является ковариантным по декларированию
% в случае наследования метода (этот случай был опущен по ошибке).
% - Уточнить символ равенства и идентичности.
%
% 2.12 - 2.13 (не было 2.11)
% - Обратить CL, где были удалены некоторые нулевые функции безопасности (чтобы включить
% публикации стабильной версии спецификации.
% - Добавить правило, что пара деклараций верхнего уровня с тем же именем базы
% - ошибка времени компиляции, за исключением случаев, когда это пара Getter/setter.
% - изменить грамматику, чтобы включить псевдонимы нефункционального типа. Правильное правило для
%, ссылаясь на «F.staticMethod()», где «F» является псевдонимом типа.
% - Добавить недостающую ошибку для циклического перенаправления генеративного конструктора.
%
% 2.8 - 2.10
% - изменить несколько предупреждений на ошибки компиляции времени, совпадающие с фактическими
% поведение инструментов.
% Устранение ошибок при коллизии названий библиотек при импорте и экспорте.
% - Уточнить спецификацию инициализации формальностей. Исправить ошибку: Добавить
% отсутствует «?» в конце для инициализации формальной функции.
% - Настройка спецификации семантики чисел JS.
% — пересмотр и уточнение разделов «Импорт и экспорт».
% - исправление ошибки: изменить <qualifiedName>, чтобы опустить один идентификатор, а затем добавить его% обратно, когда используется <qualifiedName>, как <typeIdentifier> resp. <identifier>.
% Уточнить, что ограниченный приемник функционального типа имеет метод «.call».
% - Объединить "статическую" и "инстанционную" области класса, получая один
????????????
% - Добавить спецификацию ">>" в компилируемый код JavaScript, в приложение.
% - Интеграция спецификации методов расширения в данный документ.
% - Укажите идентификационные ссылки, обозначающие члены расширения.
% - Удалить несколько нулевых функций безопасности, чтобы обеспечить стабильную публикацию
% версия спецификации, которая относится исключительно к Dart 2.10.
% - Реорганизовать спецификацию псевдонимов типа в раздел «Алиазы типов»
% (было и там, и в «генериках»).
% - Уточнить погрешность цикличности для псевдонимов типа ("F не допускается к зависимости")
% на себя.
% - Добавить ошибку для псевдонима типа <<<P000000> При создании и т.д., когда
% <<<P000001> Расширяется до переменной типа.
% - Правила правильного лексического поиска включают метод неявного расширения
% призыв.
%
% 2,7 %
% - Переименование нетерминалов "<...Определение>" - "<...Термин декларирования>" (например, это
% «классовая декларация» везде, поэтому «классовая дефиниция» непоследовательна.
% - Уточнить, что <libraryDeclaration> и <partDeclaration> являются
% - начальные символы грамматики.
% - Уточнить понятие "noSuchMethod forwarded": "m" действительно
% noSuchMethod пересылается, если реализация «m» наследуется, но
% не имеет необходимой подписи.
% - Уточнить статические проверки на "доходность" и "доходность" для обеспечения того, чтобы
Приемлемость % обеспечивается за элемент.
% - Обновить белый список для выражений типа void, чтобы включить «ожидание e»,
% в соответствии с решением по SDK #33415.
% - Повторная корректировка правил "доходности": проверки типа элемента не поддерживаются;
% заданный итерируемый/струйный должен иметь безопасный тип элемента.
% - Уточнить, что выражение типа "X расширяет T" может быть использовано, когда
% «Т» — тип функции, плюс другие подобные случаи.
% - Укажите фактический тип аргументов, передаваемых общей функции, которая вызывается
% без аргументов типа (так что это должно быть динамическое обращение).
%
% 2,6 %
% - Укажите статический анализ вызова «призывного объекта» (где
% абонент является экземпляром класса, который имеет метод «вызова».
% - Укажите, что <libraryImport> строковые буквы не могут использовать строку
% интерполяции; указать, что это динамическая ошибка для библиотеки загрузки пункт
% для загрузки библиотеки, отличной от той, которая использовалась для
Проверка % времени компиляции.
% - Укажите, что постоянное выражение "e.length" может ссылаться только на экземпляр
% геттер (в частности, он не может выполнить/сорвать расширение);
Аналогичные изменения были внесены для нескольких операторов.
% - Укажите тип аргументов полей «Вызова», полученных
% «noSuchMethod» при вызове в ответ на несостоявшийся экземпляр
% призыв.
%
% 2,4 %
% - Уточнить раздел "Экспорт".
% - Обновить правила грамматики для <declaration>, чтобы поддержать 'статический поздний финал пункт
переменные % без инициализатора; для нескольких деклараций верхнего уровня;
% для исправления существования и размещения <metadata>;
% <assignableSelectorPart>, для упрощения грамматики
% выводимых терминов; для <topLevelDeclaration>, чтобы позволить финал верхнего уровня
% и const переменные, не имеющие типа, и позволяющие верхнего уровня
% переменные, позволяющие «задержаться» в декларации переменных верхнего уровня; и
% добавление <namedParameterType> для обеспечения требуемых параметров.
% - Для поиска лексических идентификаторов используйте правила «Справочник идентификатора»
% последовательно; то есть всегда ищите «id», а также «id=» и совершайте
% к типу декларации, найденной этим поиском.
% - Укажите подпись метода «вызова» функционального объекта.
% - Добавить правило, что это ошибка для переменной типа класса.
% в нековариантном положении в суперинтерфейсе.
% - Исправьте несколько правил грамматики, в том числе: добавлено <functionExpressionBody> (для% избегать полуколона в <functionBody>, скорректированного <importSpecification>.
% - Пересмотреть раздел о каскадах. Использует композиционную грамматику и
% определяет статический тип, ошибки времени компиляции и включает «?..».
% - Исправить грамматические и лексические правила для струнных букв.
% - Изменить спецификацию «ожидать» для обеспечения безопасности по типу, избегая таких случаев, как:
% FutureOr<Future<S>>ffs = Future<S>.value(s);
% Будущее<S> fs = ожидание ffs;
%
% 2,3 %
% - Добавить требование о том, что итератор ведомости должен иметь
% тип «Итератор».
% Уточнить, какие конструкторы охвачены разделом «Постоянные»
% Конструкторов и устранение запутанной избыточности в дефинитоне
Потенциально постоянные выражения.
% - Интеграция спецификации элементов коллекции
% (UI-as-code).
%
% 2,2 %
% — Укажите, превосходят ли значения буквальных выражений Объект.
% - Разрешить объекты типа в виде выражений регистра и ключей карты const.
% - Ввести набор букв.
% - Укажите, что геттер / сеттер и метод с тем же именем базы
% ошибка, также в случае, когда класс получает оба
% суперинтерфейсы.
Укажите правило Dart 2.0, которое вы не можете реализовать, расширить или смешать
% функции.
% - обобщить спецификацию псевдонимов типа таким образом, чтобы они могли обозначать любые
Тип %, а не только функции.
% - Уточнить, что «Постоянные Конструкторы» обеспокоены неперенаправлением
Только для генеративных конструкторов.
%
% 2,1 %
% - Удалить 64-битное ограничение на целые буквы, скомпилированные в числа JavaScript.
% - позволяют целым буквам в двойном контексте оценить двойное значение.
% - Укажите динамическую ошибку для неудачного нисходящего перенаправления завода
Процентное обращение конструктора.
% - Укажите тип аргументов, принятых в перенаправляющем заводском конструкторе
Процентная декларация должна учитываться при проведении статических проверок.
% — запретить любое выражение, начинающееся с «{», а не только
% тех, которые являются буквальными.
% - Определите понятие поиска, которое необходимо для супервызовов, настройте
% спецификации супервызовов соответственно.
% Укажите, что инициализация нестатической переменной является динамической ошибкой
% с объектом, который не имеет заявленного типа (например, неисправный слепок).
% - Укажите для инициализаторов конструктора, что целевая переменная должна существовать и
% выражение инициализации должно иметь тип, который присваивается его
Тип %.
% - Укажите для суперинициализаторов, что целевой конструктор должен существовать и
% часть аргумента должна иметь соответствующую форму и тип пропуска и значение
Доля аргументов, удовлетворяющих соответствующим ограничениям.
% - Правила пересказа об абстрактных методах и наследовании для использования класса
% интерфейс.
% - Укажите, что это ошибка для конкретного класса без нетривиального
% <<<P000015> Не иметь конкретного заявления для некоторых членов
% в его интерфейсе или иметь тот, который не является правильным переопределением.
% - Использование <<<P000002> в 3 местах, исключая 2 дополнительных
% копий одного и того же текста.
% - Добавить цифру в <<<P000003> для улучшения читаемости.
% - Введите понятие поиска, которое необходимо для супервызовов.
% - Используйте новую концепцию поиска для упрощения спецификации геттера, сеттера, метода
% поиска.
% - Ввести несколько маркеров «Case<SomeTopic>» для улучшения
% читаемость.
% Реорганизовать несколько разделов, чтобы сначала указать статический анализ.
% динамическая семантика; проясните многие детали на этом пути. Секциями являются:
% <<<P000004>, <<P000005>, <<P000006>, <<P000007>,
% <<<P000008>, <<P000009>,
% <<<P000010>, <<<P000011>, <<P000012>,
% <<<P000013> и <<<P000014>.
% - Исправленная ошибка, связанная с многократным использованием одной и той же части в одном и том же
% программы с учетом экспорта.
% - Исключите все ссылки на проверенный и производственный режим, Dart 2
% не имеют режимов.
% - Интеграция спецификаций на экспедиторах noSuchMethod.
% - Укажите, что связанное удовлетворение в общих параметрах типа aliasПроцент должен подразумевать связанное удовлетворение везде в организме.
% - Укажите, что приложения псевдонимов сверхограниченного общего типа должны вызывать
Проверка правильности всех типов, происходящих в обозначенном типе.
% - Исправленный угловой вариант правил для генерации экспедиторов noSuchMethod.
% - Интеграция спецификации характеристик по параметрам, которые
% ковариантный по декларированию.
% - Интеграция спецификации характеристик по параметрам, которые
% ковариант по классу.
% - Правильный раздел «Тип функции», позволяющий вносить необходимые коррективы
% для правил, связанных с ковариантными параметрами.
% - указывает динамический тип объектов функций в нескольких контекстах, таких как:
% что можно упомянуть специальную обработку ковариантных параметров.
% — указывается, что значит для отношения переопределения быть правильным, таким образом
% добавление частей, не захваченных проверкой подтипа функции.
% - введено понятие подписей членов, указано, что они являются
% вид сущности, который содержит интерфейс класса.
% - Исправленная проверка сверхограниченности для учета дисперсии на
% верхнего уровня.
%
% 2,0 %
% - Не допускайте функции в качестве тестовых значений.
% - Начните выполнять функции «синхронно».
% - это статическое предупреждение и динамическая ошибка для назначения окончательного локального сигнала.
% — указать, что означает «эквивалент».
% - удалить @proxy.
% - Не указывайте точный объект, используемый для пустого позиционирования Аргументы и
% - Аргументы по вызову.
% - Удалить ненужное обращение с недействительными оверредами noSuchMethod.
% - Добавить >>> в качестве основного оператора.
% - Если инициализация формальной имеет тип аннотации, требуется подтип типа поля.
% - Постоянные операции "==" теперь также разрешены, если только один операнд является нулевым.
% - сделать плоский не рекурсивным.
% - Запретить реализацию двух инстанциаций одного и того же общего интерфейса.
% — обновление спецификации «FutureOr» для Dart 2.0.
% - Требовать, чтобы декларация верхнего уровня «основной» была действительной записью сценария
Декларация функции %.
% - Указать, что тип возврата сеттера или []= является недействительным, если не указано.
% - Уточнить, что "noSuchMethod" должен быть реализован, а не просто перезаявлен
% абстрактно, для устранения определенных диагностических сообщений.
% - Добавление общих функций и методов в язык.
% - Не вызывайте предупреждения, если несистемная библиотека импортирует тень системной библиотеки.
% - Обновление пересылки приложений для микширования конструкторов для правильной обработки
% необязательные параметры и конст-конструкторы.
% - Укажите «звонок» для Dart 2 (тип функции не указан для прилагаемого класса).
% - Уточнить, что ссылка на идентификатор, обозначающая верхний уровень, статический или
% местная функция оценивает клозулизацию этой декларации.
% - встраиваемые идентификаторы "mixin" и "interface".
% — сделать «асинк» *не* зарезервированным словом внутри функций асинка.
% - Добавить «Конфликты членов класса», упрощая и корректируя правила
% - конфликты, выходящие за рамки «необъявленных дважды в одном объеме».
% - указать, что целые буквы ограничены подписанными 64-битными значениями;
% и что класс «int» предназначен как подписанное 64-разрядное целое число, но
%, которые могут отличаться.
% - Укажите дисперсию и сверхсвязанные типы.
% - Ввести термины "субтерм" и "немедленный субтерм".
% - Ввести "верхний тип".
% - Конфигурируемый импорт.
% - Укажите динамический тип Итерабельного/Будущего/Потока
% вызовов функций, отмеченных sync*/async/async*.
% - Добавить приложение, в котором перечислены основные различия между 64-битными целыми числами
% и целые числа JavaScript.
% - Удалить приложение к соглашениям об именах.
% - Укажите, что «динамический» экспортируется из dart:core.
% - Удалить "булевую конверсию". Это просто ошибка - не быть дураком.
% - Настройка правила предотвращения циклических подтипов для переменных типов.
% Уточните, что использование FutureOr<T> в качестве суперинтерфейса является ошибкой.
Устраните понятие статических предупреждений, все ошибки программы теперь являются ошибками.
% — больше не является ошибкой для геттера иметь обратный тип «пустота».% - Укажите, что каждое перенаправление конструктора проверяется статически и
% динамически.
% - Укажите, что это ошибка для суперинициализатора, чтобы произойти в другом месте
%, чем в конце списка инициализаторов.
% - Обновить определения постоянной экспрессии в потенциально/время компиляции
%, так что «потенциально постоянная» зависит только от грамматики, а не от типов
% субэкспрессий.
% - сделать "==" распознать "нуль" и сделать "&&" и " | |" короткое замыкание в постоянном
% выражения.
% - Добавить выражения "as" и "is" в качестве постоянных выражений
% - Осуществлять операции '^', ' |' и '&' на постоянных операциях 'bool'.
% - Интегрировать subtyping.md. Это вводит правила Dart 2 для субтификации.
%, что, в частности, означает, что понятие «быть более конкретным»
% устраняется, а типы функций становятся контравариантными.
Типы параметров %.
% - Интегрировать инстанциацию в связанную. Это вводит понятия сырого
Типы %, отношение необработанных зависимостей и простые границы; и оно определяет
% алгоритма, используемого для расширения необработанного типа (например, «С») до
% параметризованного типа (например, «C<int>»).
% - Integrate invalid_returns.md. Это заменяет правила о том, когда
% ошибка, связанная с функцией «возврат» или «возвращение e».
% - Интегрировать generalized-void.md. Введение синтаксической поддержки для использования
% «пустота» во многих новых местах, включая аннотации переменного типа
% фактические аргументы типа; также добавляет ошибки для использования значений типа «пустота».
% - Интегрировать implicit_creation.md, указывая, как некоторые постоянные выражения
% может быть написано без "const", и все случаи "new" могут быть
% исключены.
%
% 1,15 %
% Изменить то, как языковая спецификация описывает поток управления.
% - инициализация объекта теперь правильно определяет порядок инициализации.
% - указывает, что оставление цикла ожидания должно ждать подписки
% отменяется.
% - петля ожидания только приостанавливает подписку, если она делает что-то асинхронное.
% - Утверждения Assert позволяют операнду «сообщения» и запятой.
В большинстве случаев нулевой тип в настоящее время считается подтипом всех типов.
% - Укажите, что означает NEWLINE в многострочных строках.
% — указывается тип FutureOf.
% - Утверждения могут встречаться в списках инициализаторов.
%
% 1,14 %
% - Звонок "C()", где "C" - это имя класса, является ошибкой времени компиляции.
% - Измененное описание переписываний, зависящее от функции в буквальном смысле.
% Во многих случаях переписывание не было безопасным для асинхронного кода.
% - Удаленные генерализованные разрывы.
% - Позволить "отбрасыванию" также закончить случай переключения. Разрешить брекеты вокруг корпусов переключателей.
% - Разрешить использование '=' в качестве сепаратора значений по умолчанию для названных параметров.
% Ошибка времени компиляции, если библиотека содержит одну и ту же часть дважды.
% - Теперь более конкретно о возвратных типах функций синхронизации*/синхронизации/синхронизации*
% в отношении ответных заявлений.
% - Разрешить суррогатное значение Unicode в струнных буквах.
% - сделать значение инициализации доступным в списке инициализаторов.
% - Разрешить любое выражение в утверждениях (было только условное выражение).
% - Разрешить запятые в списках аргументов и параметров.
%
% 1.11 - ECMA 408 - 4-е издание
% Укажите, что потенциально постоянные выражения должны быть действительными
%, если параметры непостоянны.
% — сделать «??» постоянным оператором времени компиляции.
Наличие нескольких неназванных библиотек больше не вызывает предупреждений.
% - Укажите операторов с нулевой осведомленностью для статических методов.
%
% 1,10
% - Разрешить смешиваниям иметь супер-классы и супер-звонки.
% - Укажите статический тип проверки для неявной переменной итератора.
% - Укажите статические типы для ряда выражений, в которых они отсутствовали.
% - делать звонки по точному типу "Функция" не вызывают предупреждения.
% - Укажите неосознанное поведение «e?.v++» и аналогичные выражения.
% - Укажите этот "пакет": URI обрабатываются в зависимости от реализации.
% - Требуется предупреждение, если форинт используется на объекте, не имеющем элемента "итератора".
%
% 1,9 - ECMA-408 - 3-е издание
%

